\documentclass[11pt]{article}

\usepackage[margin=1in]{geometry}
\usepackage{amsmath,amssymb,amsthm}
\usepackage[hidelinks]{hyperref}

\theoremstyle{plain}
\newtheorem{theorem}{Theorem}
\newtheorem{lemma}[theorem]{Lemma}
\newtheorem{proposition}[theorem]{Proposition}
\newtheorem{corollary}[theorem]{Corollary}
\theoremstyle{definition}
\newtheorem{definition}[theorem]{Definition}
\theoremstyle{remark}
\newtheorem{remark}[theorem]{Remark}

\newcommand{\conjid}{@@ID@@}

\title{@@TITLE@@\\[0.4em]\large The Last Math Competition, Conjecture \conjid}
\author{@@AUTHOR@@\\ \small @@AFFILIATION@@}
\date{@@DATE@@}

\begin{document}
\maketitle

\begin{abstract}
@@ABSTRACT@@
\end{abstract}

\section{The conjecture as stated}

The original text of conjecture \conjid{} reads:

\begin{quote}
@@STATEMENT@@
\end{quote}

\section{Reading of the statement}

% If the statement is ambiguous, under-specified, or contains an error, say so
% HERE and fix a precise reading. State explicitly which reading is being
% settled and why it is the faithful one. Do not silently repair the statement.

@@READING@@

\section{Result}

% State exactly what is established: proved, disproved, or partially resolved.

@@RESULT@@

\section{Proof}

@@PROOF@@

\section{Formalization}

The argument is formalized in Lean 4 against Mathlib in the accompanying
\texttt{lean/} directory; see \texttt{lean/Submission/Basic.lean}. The
formalization is free of \texttt{sorry} and introduces no new axioms beyond
those of Mathlib.

\section{Remarks}

% Optional: what the truth is likely to be, what stays open, why the original
% statement went wrong, what a repaired conjecture should say.

@@REMARKS@@

\end{document}
